% Statement and cited-input packet. Target proof is separately embargoed.
% Frozen II source SHA256: 74497bfe327fae25c9cdcb17334b3bf42a9793a9df7c0a4404d97568f2f760f0
% Target owner: cdp:st:C.3.11; tag: 01F6

\begin{proposition}[Proved complete cohomological dichotomy]\hypertarget{tag-01F6}{}\label{cdp:st:C.3.11}



Under the hypotheses of \ref{cdp:st:C.3.1}, a very general smooth fibre
is a two-torus, a primary or secondary Hopf surface of
algebraic dimension zero. The Inoue exclusion is \ref{cdp:st:C.3.12}.
Both Hopf
entries and the torus entry occur, by \ref{cdp:st:C.3.6} and \ref{cdp:st:D.2.7}.
For the entire smooth locus there are exactly the following
two cohomological alternatives:
\[
\begin{array}{c|c|c}
\text{smooth fibres}&
 (h^0(\OO),h^1(\OO),h^2(\OO))&
 \text{coherent direct images}\\ \hline
\text{two-tori}&(1,2,1)&
 R^2f_*\OO_Y\simeq\det R^1f_*\OO_Y\\
b_1=1,\ b_2=0&(1,1,0)&
 R^1f_*\OO_Y\simeq\OO_{\PP^1},\
 R^2f_*\OO_Y=0 .
\end{array}
\tag{C.PK.9}\hypertarget{tag-00VG}{}\label{cdp:eq:C.PK.9}
\]
The indicated cohomology dimensions hold on every
scheme-theoretic full fibre in the respective branch,
and both direct images satisfy base change there.
In the second branch
$H^q(Y,\OO_Y)=(\C,\C,0,0)$ and
$\Pic^0(Y)=\C\beta$ with the exponential convention
of \ref{cdp:st:D.1.9}. The first line records a cohomological
alternative rather than a positivity assertion about
the splitting of $R^1f_*\OO_Y$.
\end{proposition}

% Cited input cdp:st:C.3.1; statement or labelled formula only.
\begin{proposition}\hypertarget{tag-01F4}{}\label{cdp:st:C.3.1}


Let $f:Y\to C$ be a holomorphic algebraic reduction of a connected
compact complex threefold with $b_1(Y)=b_2(Y)=0$ and $a(Y)=1$.
The base is $C\simeq\PP^1$. Every smooth fibre $F$ has
\[
 c_1(K_F)=0\text{ in }H^2(F;\Q),\qquad
 c_2(F)=0\text{ in }H^4(F;\Z),\qquad
 \chi(F,\OO_F)=0,
\]
and contains no exceptional curve of the first kind. A very general
smooth fibre has $\kappa(F)\leq0$. The very general fibre's type belongs to the following
list: a complex two-torus, a bielliptic surface, a primary or secondary
Kodaira surface, a primary or secondary Hopf surface, or an Inoue
surface. Every listed type has Euler characteristic zero. This is a
preliminary necessary list. Statement \ref{cdp:st:C.3.7} below excludes the
bielliptic entry, and \ref{cdp:st:C.3.8} makes the very general Hopf fibre have
algebraic dimension zero. \ref{cdp:st:C.3.9} and \ref{cdp:st:C.3.10} exclude both Kodaira
entries. Statement \ref{cdp:st:C.3.6} realizes both primary and secondary Hopf
entries, and \ref{cdp:st:D.2.7} realizes a fixed abelian surface as general
fibre. Statement \ref{cdp:st:C.3.12} excludes the Inoue entry. Thus the
final very general list consists exactly of two-tori and
primary or secondary Hopf surfaces of algebraic dimension zero.
\end{proposition}

% Cited input cdp:st:C.3.10; statement or labelled formula only.
\begin{proposition}[Proved exclusion of primary Kodaira general fibres]\hypertarget{tag-01F5}{}\label{cdp:st:C.3.10}



A compact complex threefold with $b_1(Y)=b_2(Y)=0$ and
$a(Y)=1$ cannot have a holomorphic algebraic reduction
$f:Y\to\PP^1$ whose general fibre is a primary Kodaira
surface. Together with \ref{cdp:st:C.3.9} this excludes both Kodaira
types using only rational Betti-number vanishing.
\end{proposition}

% Cited input cdp:st:C.3.12; statement or labelled formula only.
\begin{proposition}[Proved exclusion of Inoue general fibres]\hypertarget{tag-01F7}{}\label{cdp:st:C.3.12}



Under the hypotheses of \ref{cdp:st:C.3.1}, the very general smooth fibre is
not an Inoue surface. More precisely, a proper holomorphic map from
a compact connected complex manifold to a compact Riemann surface,
whose fibres over the complement of a finite set are Inoue surfaces,
has a class in $H^1(Y;\mathbb R)$ restricting nontrivially to each
such fibre. Thus these families cannot have $b_1(Y)=0$.
This assertion permits arbitrary multiplicities and arbitrary
singularities of the omitted fibres.
\end{proposition}

% Cited input cdp:st:C.3.6; statement or labelled formula only.
\begin{proposition}[Proved actual occurrence of both Hopf types]\hypertarget{tag-01FC}{}\label{cdp:st:C.3.6}



There are compact complex threefolds with $b_1=b_2=0$ and $a=1$
whose holomorphic algebraic reduction is a smooth bundle over
$\PP^1$ with primary Hopf fibre. There are also such threefolds
with secondary Hopf fibre. In these examples every fibre has Euler
number zero and $c_3=0$.
\end{proposition}

% Cited input cdp:st:C.3.7; statement or labelled formula only.
\begin{proposition}[Proved exclusion of bielliptic fibres]\hypertarget{tag-01FD}{}\label{cdp:st:C.3.7}



A compact complex threefold $Y$ with $b_1(Y)=b_2(Y)=0$ has no
proper connected-fibre holomorphic map to a compact curve whose
general fibre is bielliptic.
\end{proposition}

% Cited input cdp:st:C.3.8; statement or labelled formula only.
\begin{proposition}[Proved exclusion of elliptic Hopf general fibres]\hypertarget{tag-01FE}{}\label{cdp:st:C.3.8}



The very general Hopf fibre of the algebraic reduction in \ref{cdp:st:C.3.1}
has algebraic dimension zero. In particular neither primary nor
secondary Hopf surfaces of algebraic dimension one can be its
general fibre. This exclusion uses no restriction-surjectivity
assertion from CDP98 or CDP20.
\end{proposition}

% Cited input cdp:st:C.3.9; statement or labelled formula only.
\begin{proposition}[Proved exclusion of secondary Kodaira general fibres]\hypertarget{tag-01FF}{}\label{cdp:st:C.3.9}



A compact complex threefold with $b_1(Y)=b_2(Y)=0$ and
$a(Y)=1$ cannot have a holomorphic algebraic reduction
$f:Y\to\PP^1$ whose general fibre is a secondary Kodaira
surface. Rational Betti-number vanishing suffices; no integral
torsion-freeness of the fibre is assumed.
\end{proposition}

% Cited input cdp:st:D.1.12; statement or labelled formula only.
\begin{proposition}\hypertarget{tag-01FX}{}\label{cdp:st:D.1.12}


For a proper connected-fibre holomorphic map $f:Y\to C$ from a
compact complex threefold with $b_2(Y)=0$ and general fibre a
two-torus or a primary Kodaira surface, the cup-product homomorphism
is an isomorphism
\[
\bigwedge^2R^1f_*\OO_Y\xrightarrow{\simeq}R^2f_*\OO_Y.
\]
Every full fibre has structure-sheaf cohomology dimensions
$(h^0,h^1,h^2)=(1,2,1)$, and these two direct images satisfy
base change at every fibre. This proves these two cases of CDP20
Proposition 5.1 by an index and cup-product argument.
\end{proposition}

% Cited input cdp:st:D.1.9; statement or labelled formula only.
\begin{proposition}\hypertarget{tag-01G7}{}\label{cdp:st:D.1.9}


Let $Y$ be a connected compact complex manifold with
$b_1(Y)=b_2(Y)=0$, and let $f:Y\to\PP^1$ be a proper
holomorphic surjection with connected fibres. There is a class
$\beta\in H^1(Y,\OO_Y)$ such that, for every irrational real number
$r$, the line bundle
\[
L_r=\exp(2\pi i\,r\beta)
\]
has infinite order on every scheme-theoretic fibre of $f$.
Every connected component of $\Pic(Y)$ likewise contains a real
one-parameter family with this property at every irrational real
parameter.
No local-freeness assumption on $R^1f_*\OO_Y$ or constant-dimension
assumption on fibre cohomology is required.
\end{proposition}

% Cited input cdp:st:D.2.7; statement or labelled formula only.
\begin{proposition}[Proved existence]\hypertarget{tag-01GE}{}\label{cdp:st:D.2.7}



There exists a smooth compact connected complex threefold $Y$
with a holomorphic map $f:Y\to\PP^1$ with connected fibres such
that
\[
 H_1(Y;\Z)=H_2(Y;\Z)=0,\qquad
 a(Y)=1,\qquad c_3(Y)=0,\qquad b_3(Y)=2.
\tag{D.21}\hypertarget{tag-0198}{}\label{cdp:eq:D.21}
\]
Every smooth fibre is the fixed abelian surface
$A=E\times E$, where $E=\C/(\Z+\mathrm i\Z)$, and the marked
integral monodromy on $H_1(A;\Z)$ is the identity. There are
exactly three singular fibres. Each is reduced with simple
normal crossings, and consists of a primary Kodaira surface
and a primary Hopf surface meeting along one smooth elliptic
curve. Consequently the trivial-monodromy class in D.2 is
nonempty, and equality in the unrestricted sign theorem
\ref{cdp:st:D.2.5} is attained.
\end{proposition}

% Cited input cdp:st:D.2.5; statement or labelled formula only.
\begin{proposition}[Proved unrestricted sign under trivial monodromy]\hypertarget{tag-01GC}{}\label{cdp:st:D.2.5}



Let $f:Y\to\PP^1$ be a proper surjective holomorphic map with
connected fibres from a compact connected complex threefold with
$b_1(Y)=b_2(Y)=0$. Suppose its smooth general fibre is a complex
two-torus and its monodromy representation on $H_1(F,\Q)$ is
trivial. Then
\[
 c_3(Y)\le0,\qquad b_3(Y)\ge2.
 \tag{D.S.1}\hypertarget{tag-019O}{}\label{cdp:eq:D.S.1}
\]
The fibres may have arbitrary singularities and arbitrary
multiplicities. In particular this proves the sign assertion of
D.2 for a holomorphic algebraic reduction; no existence assertion
is required for the conclusion.
\end{proposition}

